\documentclass[12pt]{article}
\usepackage[utf8]{inputenc}
\usepackage[T1]{fontenc}
\usepackage{hyperref}
\title{Gowro Language Resources for K-12 Curriculum Development}
\author{Compiled by Claude Sonnet 5.5 (claude-sonnet-5-5) for Andr\'as Kornai}
\date{}
\begin{document}
\maketitle

\section{Introduction}
Gowro (ISO 639-3 gwf) is a Dardic Indo-Aryan language of Indus Kohistan, Khyber Pakhtunkhwa, Pakistan, associated with the Gabar Khel clan. Speaker counts vary across sources (346 in 2019 per Wikipedia; about 1,000 in 2003). It should not be confused with Gawri (gwc), a separate language of Swat Kohistan. Searches found only encyclopedic and database entries.

\section{Primary Educational Resources}
None found.

\section{Linguistic and Reference Materials}
\begin{itemize}
\item Wikipedia, Gowro language: \url{https://en.wikipedia.org/wiki/Gowro_language}
\item Verbix Endangered Languages entry (classes Gowro as threatened): \url{https://docs.verbix.com/EndangeredLanguages/Gowro}
\end{itemize}

\section{Religious and Community Resources}
None found.

\section{Digital and Audio Resources}
Endangered Languages Project entry: \url{https://endangeredlanguages.com/lang/4255?hl=fr}

\section{Curriculum Development Context}
Shina's dominance is reported as possibly eroding Gowro use. No schooling materials or programmes were found.
\end{document}